\partdivider{B}{Nuclear Energy and Its Uses}{The binding-energy curve has a summit near iron, so heavy nuclei release energy by splitting and light ones by merging; that energy, a million times deeper than chemistry, runs reactors and stars, and the radioactive nuclei it leaves behind become clocks, tracers and medicines.}{\chref{ch:reactions} Nuclear Reactions: Fission and Fusion\\[1.2mm]\chref{ch:nucapps} Applications of Radioactivity}

\ncertchapter{4}{Nuclear Reactions: Fission and Fusion}\chlabel{ch:reactions}

So far nuclei have changed on their own. In \chref{ch:decay} an unstable nucleus emitted an $\alpha$
particle, an electron or a positron, or a gamma ray, whenever chance decided, and \chref{ch:decaylaw}
turned that chance into the exponential law. Nothing we did could hurry a decay or prevent one. This
chapter takes the active step: nuclei can also be made to change by striking them with other particles.
Once that is possible, the binding-energy curve of \chref{ch:nucleus} (Figure~\ref{fig:nuc-bea}) stops being a description and
becomes a promise. Its summit lies near iron, at about \SI{8.8}{\MeV} per nucleon, so a heavy nucleus
that splits into two middle-sized ones, or two light nuclei that merge into a heavier one, ends up more
tightly bound, and the difference in binding energy is released.

The chapter first sets up the general bookkeeping of such \emph{nuclear reactions}: what is conserved,
how much energy is released or absorbed, and how much kinetic energy a projectile must bring when energy
is absorbed. That last question is a problem in momentum conservation, and it is solved exactly as the
perfectly inelastic collisions of \chref{ch:collisions} were. The bookkeeping is then applied to the two
reactions that matter most: fission, which powers reactors, and fusion, which powers the stars. Along the
way two more pieces of Volume~I mechanics reappear: elastic collisions explain why water and graphite slow
neutrons, and the centre-of-mass frame explains why a light projectile on a heavy target wastes little
energy.

The limits deserve stating at the outset. Every energy in this chapter is computed from measured masses,
and those numbers are exact to the precision quoted. But two of the chapter's central facts, that slow
neutrons are captured far more readily than fast ones and that nuclei in the Sun fuse although they lack
the energy to touch, rest on quantum mechanics, which this book does not develop. Where that happens the
result is quoted, with a reason, and the classical estimate is shown alongside so that the size of the
gap is clear.

\section{Nuclear Reactions and Their \texorpdfstring{$Q$}{Q}-Values}\label{sec:react-q}

In 1919 Rutherford found that $\alpha$ particles passing through nitrogen gas knocked out protons, and
that the nitrogen was turned into oxygen. It was the first artificial transmutation of one element into
another:
\[
\nuc{4}{2}{He}+\nuc{14}{7}{N}\;\to\;\nuc{17}{8}{O}+\nuc{1}{1}{H}.
\]
Rutherford had no accelerator; his projectiles were the $\alpha$ particles of a radioactive source, with
kinetic energies of a few MeV. To describe such events we need a compact notation and the reaction's
energy balance. The balance is the one already used for decays in \chref{ch:decay}; the only new feature
is that there are now two particles at the start.

\begin{defbox}{Nuclear reaction and its \texorpdfstring{$Q$}{Q}-value}
In a \textbf{nuclear reaction} $\mathrm{X}(a,b)\mathrm{Y}$ a projectile $a$ strikes a target nucleus
X, usually at rest, and a light particle $b$ and a residual nucleus Y emerge:
$a+\mathrm{X}\to b+\mathrm{Y}$. Rutherford's reaction is written $\iso{14}{N}(\alpha,p)\iso{17}{O}$.
The total nucleon number $A$ and the total charge (atomic number $Z$) are the same before and after. The
\textbf{$Q$-value} is the rest energy that disappears, which equals the kinetic energy that appears:
\[
Q=\big[(m_a+m_{\mathrm X})-(m_b+m_{\mathrm Y})\big]c^2=K_b+K_{\mathrm Y}-K_a ,
\]
where the $m$ are masses and the $K$ kinetic energies (the target starts at rest). With masses in
atomic mass units, $1\,\si{\amu}\,c^2=\SI{931.494}{\MeV}$, so $Q$ in MeV is the mass difference in
\si{\amu} times $931.494$. Atomic masses may be used throughout, because total $Z$ is conserved and the
electron masses cancel. A reaction with $Q>0$ is \textbf{exoergic}: it releases energy. One with $Q<0$
is \textbf{endoergic}: the projectile must bring the missing energy in as kinetic energy.
\end{defbox}

For Rutherford's reaction, with the atomic masses $M(\iso{14}{N})=\SI{14.003074}{\amu}$,
$M(\iso{17}{O})=\SI{16.999132}{\amu}$, $M(\iso{4}{He})=\SI{4.002603}{\amu}$ and
$M(\iso{1}{H})=\SI{1.007825}{\amu}$,
\[
Q=\big[\num{14.003074}+\num{4.002603}-\num{16.999132}-\num{1.007825}\big]\times\SI{931.494}{\MeV}
=-\SI{1.19}{\MeV}.
\]
The reaction is endoergic: the products have more rest energy than the reactants. An exoergic reaction,
by contrast, can go even with a slow projectile, provided the projectile can reach the nucleus at all
(for a charged projectile the Coulomb repulsion of \chref{ch:nucleus}, Eq.~\eqref{eq:nuc-closest}, still stands in the way; we return
to this in Section~\ref{sec:react-fusion}).

\subsection{The threshold energy}

How much kinetic energy must the $\alpha$ particle bring for the endoergic reaction to go? The tempting
answer, exactly $|Q|=\SI{1.19}{\MeV}$, is wrong, and the reason is momentum. The incoming projectile
carries momentum, and no internal force can destroy it, so the products cannot be left at rest. Part of
the projectile's kinetic energy is locked into the motion of the system as a whole and can never be
turned into rest energy. \chref{ch:com} gave the tool for separating the two parts: the kinetic energy
of any system is the kinetic energy of its centre of mass (CM) plus the kinetic energy measured in the
CM frame, and only the second part can be changed by internal forces (Section~\ref{sec:cmframe}).

\begin{derivbox}{Threshold energy of an endoergic reaction}
A projectile of mass $m_a$ and kinetic energy $K_a$ strikes a target of mass $m_{\mathrm X}$ at rest.
Its momentum is $p$, with $K_a=p^2/2m_a$; the speeds are small compared with $c$, as they are whenever
$K_a$ is a few MeV and the masses are at least a nucleon's ($\approx\SI{939}{\MeV}/c^2$).
\stepn{1} \textbf{Momentum fixes the CM motion.} The total momentum $p$ is conserved, so the CM moves
at $v_{\text{cm}}=p/(m_a+m_{\mathrm X})$ before and after the reaction, and the kinetic energy of that
motion,
\[
K_{\text{cm}}=\tfrac12(m_a+m_{\mathrm X})v_{\text{cm}}^2=\frac{p^2}{2(m_a+m_{\mathrm X})}
=K_a\,\frac{m_a}{m_a+m_{\mathrm X}},
\]
is unavailable. (After the reaction the moving mass is $m_b+m_{\mathrm Y}$, which differs from
$m_a+m_{\mathrm X}$ by $|Q|/c^2$, about one part in $10^4$ here; we neglect it.)
\stepn{2} \textbf{Only the CM-frame energy can pay for $Q$.} What remains is the kinetic energy in the
CM frame,
\[
K'=K_a-K_{\text{cm}}=K_a\,\frac{m_{\mathrm X}}{m_a+m_{\mathrm X}} .
\]
The reaction can go only if $K'\ge|Q|$.
\stepn{3} \textbf{Threshold.} At threshold $K'=|Q|$ exactly, and the products are left at rest in the CM
frame, that is, moving off together at $v_{\text{cm}}$ in the laboratory. This is the perfectly
inelastic collision of \chref{ch:collisions} (Section~\ref{sec:inelastic12}), in which the energy lost to the bodies' relative
motion was also the fraction $m_{\mathrm X}/(m_a+m_{\mathrm X})$ of the projectile's kinetic energy; here
that energy becomes rest energy instead of heat. Setting $K'=|Q|$,
\begin{equation}\label{eq:react-threshold}
\boxed{K_{\text{th}}=|Q|\left(1+\frac{m_a}{m_{\mathrm X}}\right).}
\end{equation}
\stepn{4} \textbf{Check the limits.} A very heavy target ($m_{\mathrm X}\gg m_a$) barely moves, so
$K_{\text{th}}\to|Q|$: almost nothing is wasted. A projectile as heavy as the target wastes half its
energy on CM motion, and $K_{\text{th}}=2|Q|$.
\end{derivbox}

For Rutherford's reaction,
$K_{\text{th}}=1.19\,(1+\num{4.003}/\num{14.003})=\SI{1.53}{\MeV}$: about $22\%$ of the $\alpha$
particle's kinetic energy is unavoidably spent moving the whole system. The non-relativistic formula is
accurate here; the exact relativistic threshold differs from Eq.~\eqref{eq:react-threshold} by a fraction
of order $|Q|/(m_{\mathrm X}c^2)$, about $10^{-4}$.

\begin{keybox}{Threshold energy}
For an endoergic reaction ($Q<0$) on a target at rest, the projectile needs at least
$K_{\text{th}}=|Q|\,(1+m_a/m_{\mathrm X})$. A heavy target wastes little energy; a heavy projectile on a
light target wastes a great deal.
\end{keybox}

\begin{trapbox}{the threshold is not $|Q|$}
Setting $K_a=|Q|$ ignores momentum conservation: it would leave the products at rest, with zero total
momentum, although the projectile brought in momentum $p$. Nor does it matter only for exotic cases:
firing nitrogen ions at a stationary helium target, the same reaction with the roles reversed, needs
$1.19\,(1+14.003/4.003)=\SI{5.36}{\MeV}$, more than three times as much as firing $\alpha$ particles at
nitrogen.
\end{trapbox}

When $Q>0$ and the projectile is slow, the opposite question arises: how is the released energy shared?
The total momentum is then nearly zero, so the two products fly apart back to back with equal and opposite
momenta, and their kinetic energies are in inverse ratio of their masses, exactly as for the $\alpha$
particle and daughter nucleus of an $\alpha$ decay in \chref{ch:decay}. Product $b$ receives
$K_b=Q\,m_{\mathrm Y}/(m_b+m_{\mathrm Y})$.

\begin{exbox}{The first accelerator reaction}\label{ex:react-cockcroft}
In 1932 Cockcroft and Walton split lithium with protons from a high-voltage accelerator,
$\iso{7}{Li}(p,\alpha)\alpha$, that is, $p+\iso{7}{Li}\to2\,\iso{4}{He}$. Using
$M(\iso{7}{Li})=\SI{7.016003}{\amu}$, find $Q$ and the kinetic energy of each $\alpha$ particle when the
proton is very slow.
\stepn{1} \textbf{$Q$ from atomic masses.}
$Q=[\num{7.016003}+\num{1.007825}-2(\num{4.002603})]\times\SI{931.494}{\MeV}=\SI{17.3}{\MeV}$, exoergic.
\stepn{2} \textbf{Sharing.} With negligible initial momentum the two $\alpha$ particles fly apart back to
back. They have equal masses, so they share equally: \SI{8.7}{\MeV} each.
\stepn{3} \textbf{Meaning.} The protons were accelerated through only a few hundred kilovolts, yet each
reaction released \SI{17.3}{\MeV}. Cockcroft and Walton measured the $\alpha$ energies and found them in
agreement with the mass difference: the first direct laboratory check of $E=mc^2$.
\ans{$Q=\SI{17.3}{\MeV}$; \SI{8.7}{\MeV} each}
\end{exbox}

Reactions are also how radioactive isotopes are manufactured. \textbf{Neutron capture},
$\mathrm X(n,\gamma)\mathrm Y$, is the workhorse in a reactor: the neutron feels no Coulomb barrier, and
the binding energy it gains on joining the nucleus (several MeV) leaves as a gamma ray. Cobalt-60, made
from ordinary cobalt-59 in this way, was the example of production at a constant rate in
\chref{ch:decaylaw} (Section~\ref{sec:decaylaw-activation}), and most medical isotopes of \chref{ch:nucapps} are made by neutron capture or by
reactions in accelerators. Reactions also revealed the neutron itself: in 1932 Chadwick identified the
penetrating radiation from $\iso{9}{Be}(\alpha,n)\iso{12}{C}$ as uncharged particles of almost the
proton's mass. The most important reactions of all, however, are those in which a nucleus splits or two
nuclei merge, and the binding-energy curve says which way each must go.

\section{Splitting the Nucleus: Fission}\label{sec:react-fission}

In December 1938 Otto Hahn and Fritz Strassmann bombarded uranium with neutrons and found barium among
the products, a nucleus about half as heavy. Lise Meitner and Otto Frisch explained it within weeks using
the liquid-drop picture of \chref{ch:nucleus}. The absorbed neutron leaves the uranium ``drop'' excited
and oscillating. If it stretches into a dumbbell, the short-range nuclear attraction, which acts only
between neighbouring nucleons, can no longer hold the two lobes, while the long-range Coulomb repulsion
between all the protons of one lobe and all those of the other keeps pushing, and the drop splits
(Figure~\ref{fig:react-fission}). They called it \textbf{fission}, after the division of a living cell.

\begin{figure}[htbp]\centering
\tikzfig{react_fission}
\figcaption{fig:react-fission}{Fission in the liquid-drop picture. Absorbing a slow neutron leaves
uranium-236 excited and oscillating; once it stretches enough to form a neck, the short-range nuclear
attraction can no longer hold the two halves against their long-range electric repulsion. The fragments
fly apart back to back with equal and opposite momenta, so the lighter (krypton) is the faster, along
with two or three neutrons.}
\end{figure}

A typical outcome is
\[
\nuc{1}{0}{n}+\nuc{235}{92}{U}\;\to\;\nuc{236}{92}{U}^{*}\;\to\;\nuc{141}{56}{Ba}+\nuc{92}{36}{Kr}+3\,\nuc{1}{0}{n},
\]
where the asterisk marks an excited nucleus. Check the bookkeeping: $1+235=141+92+3=236$ nucleons and
$92=56+36$ protons. The split is not unique; uranium-236 breaks in dozens of ways, most often into one
fragment near $A\approx95$ and one near $A\approx140$, with two or three neutrons on average.

Why are \emph{neutrons} the trigger, rather than protons or $\alpha$ particles? A charged projectile must
first climb the Coulomb barrier of a $Z=92$ nucleus. With $R=R_0A^{1/3}$ and $R_0=\SI{1.2}{\fm}$ for both
partners and $e^2/4\pi\varepsilon_0=\SI{1.44}{\MeV\fm}$ (both from \chref{ch:nucleus}), a proton
touching uranium-235 has potential energy $92(1.44)/[1.2(1+235^{1/3})]\approx\SI{15}{\MeV}$, and an
$\alpha$ particle about \SI{28}{\MeV}. A neutron carries no charge, feels no repulsion, and can drift in
at any speed, however slow. That is why Chadwick's discovery of the neutron in 1932 opened the way to
fission six years later.

\subsection{The energy released}

With the atomic masses $M(\iso{235}{U})=\SI{235.043930}{\amu}$, $M(\iso{141}{Ba})=\SI{140.914403}{\amu}$,
$M(\iso{92}{Kr})=\SI{91.926173}{\amu}$ and the neutron mass $m_n=\SI{1.008665}{\amu}$,
\[
Q=\big[\num{235.043930}+\num{1.008665}-\num{140.914403}-\num{91.926173}-3(\num{1.008665})\big]\times\SI{931.494}{\MeV}
=\SI{173}{\MeV}.
\]
The binding-energy curve says the same thing without a mass table, and says \emph{why}. Since every
nucleus is lighter than its separated nucleons by its binding energy $\Eb$, the rest energy that
disappears in a reaction is the binding energy gained: $Q$ equals the total binding energy of the
products minus that of the reactants (free neutrons have none), which is Eq.~\eqref{eq:nuc-qb} of \chref{ch:nucleus}. Example~\ref{ex:nuc-fission} there made the
estimate: a uranium nucleus has $\Eb/A\approx\SI{7.6}{\MeV}$, nuclei near $A\approx120$ have about
\SI{8.5}{\MeV}, and raising about $236$ nucleons from the first value to the second releases about
$236\times\SI{0.9}{\MeV}\approx\SI{200}{\MeV}$.

\begin{keybox}{Fission and fusion on the binding-energy curve}
By Eq.~\eqref{eq:nuc-qb}, $Q$ is the total binding energy gained, so moving nucleons up the
binding-energy curve releases energy. Heavy nuclei move up by splitting (about
\SI{0.9}{\MeV} per nucleon for uranium); light nuclei move up, much more steeply, by merging.
\end{keybox}

The fragments are themselves very neutron-rich. Heavy stable nuclei have a neutron-to-proton ratio
$(A-Z)/Z\approx1.55$, middle-mass stable nuclei only about $1.3$, so fragments born with uranium's
proportions lie well above the valley of stability of \chref{ch:decay}. Each undergoes a chain of
$\beta^-$ decays towards the valley, and these decays (with their gamma rays) bring the total release per
fission to about \SI{200}{\MeV}.

\begin{trapbox}{one split's \SI{173}{\MeV} is not the \SI{200}{\MeV} average}
The $Q$-value of one particular split (here \SI{173}{\MeV}) counts only the prompt step. The figure of
about \SI{200}{\MeV} is an average over the many possible splits, \emph{including} the later $\beta$ and
$\gamma$ decays of the fragments. Use masses when a specific reaction is given, and \SI{200}{\MeV} for
reactor and fuel estimates, unless the problem states otherwise.
\end{trapbox}

Where does the energy go? Table~\ref{tab:react-fission-energy} gives the average shares. Almost all of it appears
first as kinetic energy of the two fragments, driven apart by their Coulomb repulsion, and is turned into
heat as they plough through the surrounding fuel within a fraction of a millimetre.

\begin{tablebox}{Where the energy of a uranium-235 fission goes (average, approximate)}
\centering
\begin{tabular}{@{}lr@{}}\toprule
Carrier & Energy (MeV)\\\midrule
kinetic energy of the two fragments & $170$\\
kinetic energy of the prompt neutrons & $5$\\
prompt gamma rays & $7$\\
later $\beta$ and $\gamma$ decays of the fragments & $13$\\
antineutrinos from those $\beta$ decays (escape) & $10$\\\midrule
total & $\approx205$\\\bottomrule
\end{tabular}\label{tab:react-fission-energy}
\par\raggedright
The antineutrinos pass through the reactor and the Earth and are lost. In a reactor this loss is roughly
made up by gamma rays emitted when the surplus neutrons are captured in the fuel and structure, so the heat
deposited is close to \SI{200}{\MeV} per fission, the round figure used for all estimates.
\end{tablebox}

\begin{remarkbox}{Why uranium-235 and not uranium-238?}
Absorbing a neutron releases its binding energy in the new nucleus as internal excitation. Nucleons are
more tightly bound in pairs (the pairing term of the liquid-drop formula in \chref{ch:nucleus}). A neutron
that completes a pair, turning uranium-235 (odd neutron number $A-Z=143$) into uranium-236 (even, $144$),
releases about \SI{6.5}{\MeV}, more than the roughly \SI{6}{\MeV} needed to deform the nucleus to the
point of splitting (the \emph{fission barrier}): even a nearly stationary neutron triggers fission. A
neutron absorbed by uranium-238 (even, $146$) is left unpaired in uranium-239 and releases only about
\SI{4.8}{\MeV}, too little. Uranium-238 therefore fissions only with fast neutrons, above about
\SI{1}{\MeV}, whose kinetic energy makes up the difference. Natural uranium is $99.3\%$ uranium-238 and
only $0.72\%$ uranium-235, which is the root of the engineering problems of the next two sections.
\end{remarkbox}

To appreciate \SI{200}{\MeV}, compare it with chemistry: burning one molecule of methane releases about
\SI{9}{eV}. The ratio per reaction is twenty million. Per kilogram of fuel the comparison is just as
striking.

\begin{exbox}{Energy per kilogram}\label{ex:react-perkg}
Taking \SI{200}{\MeV} per fission, find the energy released by complete fission of \SI{1.00}{kg} of
uranium-235, and compare with coal, whose heat of combustion is about \SI{30}{MJ/kg}.
\stepn{1} \textbf{Count nuclei.} One mole is \SI{235.0}{g}, so
$N=(1000/235.0)(\num{6.022e23})=\num{2.56e24}$ nuclei.
\stepn{2} \textbf{Energy.} $E=\num{2.56e24}\times200\times\SI{1.602e-13}{J}=\SI{8.2e13}{J}$.
\stepn{3} \textbf{Compare.} $\num{8.2e13}/\num{3.0e7}=\num{2.7e6}$: a kilogram of uranium-235 does the
work of about $2700$ tonnes of coal.
\ans{\SI{8.2e13}{J}, about $2.7$ million times coal}
\end{exbox}

\section{The Chain Reaction}\label{sec:react-chain}

A single fission is a laboratory curiosity; what makes fission an energy source is that each one releases
two or three neutrons (2.4 on average for uranium-235), and each of those can cause another fission. The
process can feed itself. Whether it grows, holds steady or dies out depends on a single number.

\begin{defbox}{Multiplication factor}
The \textbf{multiplication factor} $k$ is the average number of neutrons from one fission that go on to
cause another fission. (In this volume $k$ always means this; in Volume~I it was a spring constant.) It is
smaller than the $2.4$ neutrons released, because some neutrons leak out of the fuel and some are captured
without causing fission. If one generation contains $n_0$ fissions, generation $n$ contains $n_0k^n$:
\begin{itemize}[itemsep=1pt]
\item $k<1$: \textbf{subcritical}, the chain dies out;
\item $k=1$: \textbf{critical}, a steady chain (a reactor at constant power);
\item $k>1$: \textbf{supercritical}, exponential growth (a reactor raising its power, or a bomb).
\end{itemize}
\end{defbox}

Figure~\ref{fig:react-chain} shows the growth for $k=3$, an idealisation with no losses. The growth is
geometric, generation by generation, and it is fast: with $k=2$, eighty generations multiply the
population by $2^{80}\approx10^{24}$.

\begin{figure}[htbp]\centering
\tikzfig{react_chain}
\figcaption{fig:react-chain}{A chain reaction with $k=3$: each fission's neutrons cause three more
fissions, so the generations hold $1$, $3$, $9$, \dots, $3^n$ fissions. In practice some neutrons escape
(dashed) or are absorbed without fission, and a reactor is held precisely at $k=1$.}
\end{figure}

A subcritical assembly fed by an external neutron source is also instructive. One source neutron
causes on average $k$ neutrons in the next generation, $k^2$ in the one after, and so on, so the total
number of neutrons it gives rise to, counting itself, is the geometric series
\[
1+k+k^2+\dots=\frac{1}{1-k}\qquad(k<1),
\]
which follows from multiplying the sum $S$ by $k$: $S-kS=1$. With $k=0.9$ each source neutron is
multiplied tenfold. As $k\to1$ the sum diverges: the chain no longer needs the source. This is how a
reactor approaching criticality is watched; the count from a neutron detector grows as $1/(1-k)$ and
warns long before $k=1$ is reached.

\begin{keybox}{Critical mass: surface against volume}
Neutrons are produced throughout the volume of the fuel ($\propto R^3$ for a sphere of radius $R$) but
escape through its surface ($\propto R^2$), so the fraction that escapes falls roughly as $1/R$. Below a
certain size too many neutrons leak out and $k<1$, however pure the fuel. The smallest mass that can
sustain the chain is the \textbf{critical mass}, about \SI{50}{kg} for a bare sphere of pure
uranium-235. A neutron-reflecting shell around the fuel, or compressing it, reduces the critical mass.
Fissile material is safe when divided into pieces each well below critical.
\end{keybox}

\section{Taming the Chain: The Reactor}\label{sec:react-reactor}

A bomb lets $k$ exceed $1$ by as much as possible; a reactor holds $k=1$ to within a fraction of a
percent, for years, and extracts the heat. Figure~\ref{fig:react-reactor} shows the most common design,
the pressurised-water reactor. Its parts each answer one physical problem.

\begin{figure}[htbp]\centering
\tikzfig{react_reactor}
\figcaption{fig:react-reactor}{Schematic pressurised-water reactor (not to scale). The water in the
vessel is both moderator and coolant; the control rods set $k$; the hot water (orange) heats a separate
loop in the steam generator and returns cool (blue), and the steam drives a turbine and generator. Only
about a third of the heat becomes electricity.}
\end{figure}

\begin{itemize}
\item \textbf{Fuel.} Uranium oxide pellets sealed in metal rods. Natural uranium is only $0.72\%$
uranium-235. The rest, uranium-238, captures many neutrons without fissioning, so most reactors use
\emph{enriched} uranium, in which the uranium-235 fraction has been raised to $3$--$5\%$ by physically
separating the isotopes. (Gas centrifuges spin uranium hexafluoride, and the slightly lighter
uranium-235 molecules concentrate near the axis.) Heavy-water reactors can run on natural uranium; weapons
require about $90\%$.
\item \textbf{Moderator.} Fission neutrons are fast, about \SI{2}{\MeV}, but uranium-235 is hundreds of
times more likely to absorb a slow \emph{thermal} neutron, one whose kinetic energy is about
\SI{0.025}{eV}, the energy of thermal motion at room temperature. The moderator, a material of light
nuclei, slows the neutrons by elastic collisions: ordinary water, heavy water ($\mathrm{D_2O}$) or
graphite.
\item \textbf{Control rods.} Cadmium or boron absorbs neutrons strongly. Inserting the rods lowers $k$;
withdrawing them raises it.
\item \textbf{Coolant.} Carries the heat to a steam generator, where steam is raised to drive a turbine.
In a pressurised-water reactor the moderator water is also the coolant.
\item \textbf{Shielding.} Thick concrete and steel absorb the neutrons and gamma rays.
\end{itemize}

Why should a slow neutron be captured more readily than a fast one? A rough classical picture is that a
slow neutron lingers. A thermal neutron moves at $v=\sqrt{2K/m_n}\approx\SI{2200}{m/s}$, while a
\SI{2}{\MeV} neutron moves at about \SI{2.0e7}{m/s}, roughly $9000$ times faster; the slow one spends
$9000$ times longer near each nucleus it passes. The proper explanation is quantum mechanical and is
quoted here rather than derived: every particle also behaves as a wave of \emph{de~Broglie wavelength}
$h/p$, with $h$ Planck's constant and $p$ the momentum, and the slow neutron's wavelength is about $9000$
times longer, so it interacts with a nucleus over a far larger effective area. The measured fission
probability of uranium-235 is indeed about $500$ times larger for thermal neutrons than for
\SI{2}{\MeV} ones.

\subsection{Moderation by elastic collisions}

The moderator must remove almost all of a neutron's kinetic energy, from \SI{2}{\MeV} to
\SI{0.025}{eV}, a factor of nearly $10^8$, without absorbing the neutron. Collisions at these energies
leave both nuclei in their ground states, so kinetic energy is conserved: they are elastic collisions, and
\chref{ch:collisions} already solved them.

\begin{derivbox}{Why light nuclei make good moderators}
A neutron of mass $m$ and velocity $v_0$ hits a stationary nucleus of mass $Am$ head-on and elastically.
(The nucleus's mass is very nearly $A$ neutron masses, close enough here.) Let the neutron's velocity
afterwards be $v_1$ and the nucleus's $V$, all along one line, signed positive in the direction of $v_0$.
\stepn{1} \textbf{The two conservation laws.}
\begin{align*}
\text{momentum:}&\quad mv_0=mv_1+AmV &&\Rightarrow\quad v_0-v_1=AV,\\
\text{kinetic energy:}&\quad \tfrac12mv_0^2=\tfrac12mv_1^2+\tfrac12AmV^2 &&\Rightarrow\quad (v_0-v_1)(v_0+v_1)=AV^2 .
\end{align*}
\stepn{2} \textbf{Divide.} Dividing the energy equation by the momentum equation gives $v_0+v_1=V$: the
relative velocity is reversed, as for every elastic collision in \chref{ch:collisions}.
\stepn{3} \textbf{Solve.} Substituting this $V$ into $v_0-v_1=AV$,
\[
v_0-v_1=A(v_0+v_1)\quad\Longrightarrow\quad v_1=-\frac{A-1}{A+1}\,v_0 ,
\]
which is Volume~I's $v_1=\frac{m_1-m_2}{m_1+m_2}u_1$ with $m_1=m$, $m_2=Am$. The minus sign means the
neutron bounces back if $A>1$.
\stepn{4} \textbf{Energy kept.} The fraction $\phi$ of its kinetic energy that the neutron \emph{keeps}
is $v_1^2/v_0^2$:
\begin{equation}\label{eq:react-phi}
\boxed{\phi=\frac{K_1}{K_0}=\left(\frac{A-1}{A+1}\right)^2 .}
\end{equation}
\stepn{5} \textbf{Limits.} For hydrogen ($A=1$), $\phi=0$: a head-on collision stops the neutron dead,
the equal-mass exchange of Volume~I's Newton's cradle. For carbon ($A=12$) the neutron keeps
$(11/13)^2=71.6\%$. For uranium ($A=238$) it keeps $(237/239)^2=98.3\%$ and barely slows, like a
ping-pong ball bouncing off a bowling ball.
\end{derivbox}

\begin{keybox}{Energy kept per head-on elastic collision}
A neutron striking a nucleus of mass number $A$ at rest, head-on and elastically, keeps the fraction
$\phi=\big((A-1)/(A+1)\big)^2$ of its kinetic energy. After $n$ such collisions $K_n=K_0\phi^n$, so the
number needed to go from $K_0$ to $K_n$ is $n=\ln(K_n/K_0)/\ln\phi$, rounded up.
\end{keybox}

Because $K_n/K_0=\phi^n$, the logarithm of the energy falls by the same step at every collision, and on a
logarithmic scale each moderator traces a straight line (Figure~\ref{fig:react-moderation}). The slopes
make the ranking plain: the lighter the nucleus, the steeper the fall. Solving
$\phi^n=K_n/K_0$ for $n$ takes one logarithm (see \chref{ch:toolkit}, Section~\ref{sec:explog}).

\begin{figure}[htbp]\centering
\tikzfig{react_moderation}
\figcaption{fig:react-moderation}{Kinetic energy of a neutron after $n$ head-on elastic collisions,
on a logarithmic scale, for three moderating nuclei (colours distinguish the materials). Each collision
multiplies the energy by the same $\phi$, so each line is straight. Deuterium reaches the thermal level
in $9$ collisions and carbon in $55$; uranium-238 would need about $1100$.}
\end{figure}

\begin{exbox}{Graphite as a moderator}\label{ex:react-graphite}
How many head-on collisions with carbon-12 nuclei slow a \SI{2.0}{\MeV} neutron to \SI{0.025}{eV}?
\stepn{1} \textbf{Energy kept per collision.} $\phi=(11/13)^2=0.716$.
\stepn{2} \textbf{Number of collisions.}
$n=\dfrac{\ln(0.025/\num{2.0e6})}{\ln 0.716}=\dfrac{-18.2}{-0.334}=54.5$, so $55$ collisions.
\stepn{3} \textbf{Reality check.} Real collisions are mostly glancing, as the next remark explains, and
about $115$ are needed in graphite. The head-on count is a lower bound.
\ans{$55$ (head-on); about $115$ in practice}
\end{exbox}

\begin{remarkbox}{Glancing collisions}
A head-on collision is the most effective one; a glancing collision deflects the neutron and takes less
energy. Averaging over all directions of scattering (which needs the two-dimensional collisions of
\chref{ch:collisions} and a result about the distribution of directions that we state without proof)
gives the average loss of $\ln K$ per collision, $\xi=1+\dfrac{\phi\ln\phi}{1-\phi}$, with $\xi=1$ for
hydrogen. The number of collisions to thermalise a \SI{2}{\MeV} neutron is then
$\ln(\num{2e6}/0.025)/\xi$: about $18$ in hydrogen, $25$ in deuterium and $115$ in carbon. The ranking is
the head-on ranking; only the counts grow.
\end{remarkbox}

Ordinary hydrogen is the best slower, but it is not the best moderator, because it also captures
neutrons: $n+p\to\iso{2}{H}+\gamma$. A good moderator must slow neutrons \emph{and} leave them alive.

\begin{remarkbox}{Heavy water and India's three-stage programme}
Deuterium slows neutrons almost as well as hydrogen and hardly absorbs them, so heavy-water reactors can
run on \emph{natural} uranium, while light-water reactors need enriched fuel. India's pressurised
heavy-water reactors are the first stage of a three-stage plan designed around the country's modest
uranium and very large thorium reserves. Stage one burns natural uranium and breeds plutonium-239 from
uranium-238 by neutron capture. Stage two uses that plutonium in fast breeder reactors. Stage three
breeds fissile uranium-233 from thorium-232:
\[
\iso{232}{Th}+n\to\iso{233}{Th}\xrightarrow{\ \beta^-\ }\iso{233}{Pa}\xrightarrow{\ \beta^-\ }\iso{233}{U}.
\]
\end{remarkbox}

\subsection{Delayed neutrons and control}

There is a puzzle in holding $k=1$. In a thermal reactor the time from one fission to the next, the
\emph{generation time}, is only about \SI{e-4}{s}. If $k$ drifted to $1.001$, the power would grow by
$0.1\%$ every \SI{e-4}{s}, a factor of $e$ in a tenth of a second; no mechanical control rod could keep up.
Reactors are controllable for a reason that is easy to overlook.

\begin{keybox}{Delayed neutrons make reactors controllable}
About $0.65\%$ of uranium-235 fission neutrons are emitted not at once but seconds later (about
\SI{13}{s} on average), by certain $\beta$-decaying fragments. Averaged over all neutrons, the generation
time becomes about $(0.9935)(\SI{e-4}{s})+(0.0065)(\SI{13}{s})\approx\SI{0.1}{s}$. A reactor that is
critical only with the help of its delayed neutrons therefore changes power over tens of seconds instead
of milliseconds, slowly enough for control rods to respond.
\end{keybox}

\begin{exbox}{A reactor slightly supercritical}\label{ex:react-supercritical}
A reactor is made slightly supercritical, with $k=1.001$. How many generations does it take for the power
to rise $100$-fold? How long is that with prompt neutrons alone (generation time \SI{e-4}{s}), and with
the delayed neutrons included (effective generation time \SI{0.1}{s})?
\stepn{1} \textbf{Generations.} $k^n=100$ gives $n=\ln100/\ln1.001=4607$, about $4600$ generations.
\stepn{2} \textbf{Times.} Prompt only: $4600\times\SI{e-4}{s}=\SI{0.46}{s}$. With delayed neutrons:
$4600\times\SI{0.1}{s}=\SI{460}{s}$, nearly eight minutes.
\stepn{3} \textbf{Meaning.} Half a second is an uncontrollable excursion; eight minutes is routine work
for the control system.
\ans{$4600$ generations; \SI{0.46}{s} or about \SI{8}{min}}
\end{exbox}

\begin{trapbox}{delayed neutrons help only while $k-1$ is smaller than $0.0065$}
The delayed neutrons slow the growth only if the chain \emph{needs} them. If $k$ exceeds $1$ by more than
the delayed fraction, the prompt neutrons alone give a multiplication $k(1-0.0065)>1$ and the power grows
on the prompt time scale whatever the delayed neutrons do. For $k=1.01$, for instance,
$k(1-0.0065)=1.0034$. This condition, \emph{prompt criticality}, is what every reactor design must exclude;
the Chernobyl reactor reached it in 1986 and its power rose many times over within seconds.
\end{trapbox}

Finally, how much fuel does a reactor burn? Because each fission releases so much, astonishingly little.

\begin{exbox}{Fuel consumption of a power reactor}\label{ex:react-fuel}
A reactor produces \SI{1000}{MW} of thermal power. Taking \SI{200}{\MeV} per fission, find the fission
rate and the mass of uranium-235 consumed per day.
\stepn{1} \textbf{Energy per fission.} $200\times\SI{1.602e-13}{J}=\SI{3.20e-11}{J}$.
\stepn{2} \textbf{Fission rate.} $\num{1.000e9}/\num{3.20e-11}=\SI{3.12e19}{s^{-1}}$.
\stepn{3} \textbf{Mass per day.} In \SI{86400}{s}, $\num{3.12e19}\times\num{86400}=\num{2.70e24}$ nuclei,
that is, $(\num{2.70e24}/\num{6.022e23})\times\SI{235.0}{g}=\SI{1.05}{kg}$.
\ans{\SI{3.12e19}{s^{-1}}; \SI{1.05}{kg} per day}
\end{exbox}

A thousand megawatts of heat is not a thousand megawatts of electricity. The heat is turned into work by a
steam turbine, a heat engine, and \chref{ch:thermo} showed that no heat engine working between a hot and a
cold reservoir can convert all its heat into work. With steam at about \SI{600}{K} and cooling water near
\SI{300}{K}, the Carnot limit is $1-300/600=50\%$; real plants achieve about one third. The rest of the
heat leaves through the cooling towers, as it does from a coal-fired plant.

\section{Joining Nuclei: Fusion}\label{sec:react-fusion}

On the light side of the binding-energy curve (Figure~\ref{fig:nuc-bea} and Section~\ref{sec:nuc-curve}), combining nuclei releases energy, and
per nucleon far more than fission does. Deuterium--tritium fusion releases \SI{17.6}{\MeV} among $5$
nucleons, about \SI{3.5}{\MeV} per nucleon, against $200/236\approx\SI{0.85}{\MeV}$ per nucleon for
uranium fission. The curve rises much more steeply on the light side than it falls on the heavy side. The
obstacle is electrical: two nuclei must approach within the range of the strong force, a few
femtometres, against their Coulomb repulsion (Figure~\ref{fig:react-barrier}).

\begin{figure}[htbp]\centering
\tikzfig{react_barrier}
\figcaption{fig:react-barrier}{Potential energy of two nuclei against the separation of their centres
(schematic, not to scale). Outside contact the Coulomb repulsion rises as $1/r$ to the barrier height
$U_B$ at $r=R_1+R_2$; inside, the strong force takes over. A nucleus with the thermal energy shown would
be turned back classically at the turning point. In the Sun's core the typical thermal energy is several
hundred times below $U_B$, and fusion proceeds by quantum tunnelling through the barrier.}
\end{figure}

\begin{derivbox}{The Coulomb barrier and the temperature it seems to require}
\stepn{1} \textbf{Barrier height.} Two protons touching at a centre-to-centre distance $r\approx\SI{2}{\fm}$
have electric potential energy
\[
U_B=\frac{e^2}{4\pi\varepsilon_0\,r}=\frac{\SI{1.44}{\MeV\fm}}{\SI{2}{\fm}}=\SI{0.72}{\MeV}.
\]
For nuclei of charges $Z_1e$ and $Z_2e$ and radii $R_1$, $R_2$, $U_B=Z_1Z_2e^2/[4\pi\varepsilon_0(R_1+R_2)]$.
\stepn{2} \textbf{Temperature.} In a gas at absolute temperature $T$ the average translational kinetic
energy of a particle is $\tfrac32k_BT$ (\chref{ch:kinetic}; $k_B=\SI{1.381e-23}{J/K}$ is Boltzmann's
constant). If this average had to equal the barrier,
\[
T=\frac{2U_B}{3k_B}=\frac{2(0.72\times\SI{1.602e-13}{J})}{3(\SI{1.381e-23}{J/K})}\approx\SI{5.6e9}{K}.
\]
\stepn{3} \textbf{Compare with the Sun.} The Sun's core is at only \SI{1.5e7}{K}, about $370$ times
colder, where $k_BT\approx\SI{1.3}{\keV}$, some $560$ times below the barrier.
\end{derivbox}

Yet the Sun shines. Two effects close the gap, and the second is beyond classical physics. (i) Particle
speeds in a gas are spread out (the Maxwell distribution of \chref{ch:kinetic}), and a small fraction of
the protons in its high-energy tail move several times faster than average. (ii) Quantum mechanics lets a
particle pass through a barrier it cannot climb, a process called \textbf{tunnelling}, with a probability
that falls off exponentially with the barrier's height and width and with the particle's mass. We quote
tunnelling rather than derive it, because the derivation needs the wave mechanics of a later course; all
that is needed here is that it exists and that it is very improbable for a single encounter. The
product of the two small factors is tiny, which is exactly why the Sun burns slowly and steadily for ten
billion years instead of exploding.

\begin{trapbox}{fusion does not need $\tfrac32k_BT=U_B$}
The estimate $T=2U_B/3k_B$ is the temperature at which an \emph{average} nucleus could climb the barrier
classically. Real fusion runs hundreds of times colder, on the rare fast nuclei of the Maxwell tail and on
tunnelling. Quote the classical estimate only as an upper bound that shows why fusion is hard.
\end{trapbox}

\begin{exbox}{The deuterium--tritium barrier}\label{ex:react-dtbarrier}
Estimate the Coulomb barrier between a deuteron and a triton, taking their radii from
$R=R_0A^{1/3}$ with $R_0=\SI{1.2}{\fm}$, and the temperature at which $\tfrac32k_BT$ equals it.
\stepn{1} \textbf{Contact distance.} $r=1.2(2^{1/3}+3^{1/3})=1.2(1.260+1.442)=\SI{3.24}{\fm}$.
\stepn{2} \textbf{Barrier.} Both charges are $+e$: $U_B=1.44/3.24=\SI{0.444}{\MeV}$.
\stepn{3} \textbf{Classical temperature.} $T=2U_B/3k_B=\SI{3.4e9}{K}$. Practical D--T plasmas work at
about \SI{e8}{K}, thirty times colder, relying on the fast tail and on tunnelling.
\ans{$U_B\approx\SI{0.44}{\MeV}$; $T\approx\SI{3.4e9}{K}$}
\end{exbox}

\section{How the Sun Shines: The Proton--Proton Chain}\label{sec:react-sun}

The Sun is about three-quarters hydrogen by mass, and in its core hydrogen is fused into helium by a sequence of
reactions called the proton--proton chain (Figure~\ref{fig:react-ppchain}). Its main branch runs
\begin{align*}
&\textbf{1}\quad p+p\to\iso{2}{H}+e^++\nu_e &&\text{(weak interaction; extremely slow)}\\
&\textbf{2}\quad \iso{2}{H}+p\to\iso{3}{He}+\gamma\\
&\textbf{3}\quad \iso{3}{He}+\iso{3}{He}\to\iso{4}{He}+2p
\end{align*}
with steps 1 and 2 occurring twice for each helium nucleus made. Here $e^+$ is the positron and $\nu_e$ the
neutrino of the $\beta^+$ decays of \chref{ch:decay}. The net result is
\[
4p\;\to\;\iso{4}{He}+2e^++2\nu_e .
\]

\begin{figure}[htbp]\centering
\tikzfig{react_ppchain}
\figcaption{fig:react-ppchain}{The main branch of the proton--proton chain. Steps 1 and 2 occur twice for
each helium-4 made, and step 3 returns two of the six protons. Step 1 needs a proton to turn into a neutron
by the weak interaction during a brief encounter; a given proton in the Sun's core waits on average about
ten billion years for this, which is what makes the Sun's fuel last.}
\end{figure}

Using atomic masses for the net reaction,
\[
Q=\big[4M(\iso{1}{H})-M(\iso{4}{He})\big]c^2=\big[4(\num{1.007825})-\num{4.002603}\big]\times\SI{931.494}{\MeV}=\SI{26.7}{\MeV}.
\]
The electron bookkeeping deserves a moment, because here, unlike in fission, the charge changes hands. The
four hydrogen atoms on the left include four electrons; the helium atom on the right has two. The two
``extra'' electrons annihilate with the two positrons, and that annihilation energy ($4m_ec^2$) is
automatically counted when atomic masses are used, as it was for $\beta^+$ decay in \chref{ch:decay}.
About \SI{0.5}{\MeV} of the \SI{26.7}{\MeV} is carried off by the two neutrinos, which escape the Sun
unhindered; the rest heats the core.

\begin{remarkbox}{The Sun loses four million tonnes a second}
The Sun radiates a power $P_\odot=\SI{3.83e26}{W}$. By $E=mc^2$ it loses mass at
$P_\odot/c^2=\SI{4.3e9}{kg/s}$, over four million tonnes every second. Since \SI{26.7}{\MeV} is released per
four hydrogen atoms, it converts about $\SI{6.0e11}{kg}$ of hydrogen into helium every second. Even so,
over its $4.6$ billion years it has used only about $6\%$ of its hydrogen at today's rate (less in fact,
since the young Sun was fainter): its mass is \SI{2.0e30}{kg}, about $71\%$ of it hydrogen.
\end{remarkbox}

\subsection{Fusion on Earth}

A reactor on Earth cannot wait billions of years for the weak interaction, nor can it use a star's
enormous mass to hold its fuel together. Terrestrial fusion aims at the easiest reaction,
deuterium--tritium:
\[
\iso{2}{H}+\iso{3}{H}\to\iso{4}{He}+n,\qquad Q=\SI{17.6}{\MeV}.
\]
It is the easiest for three reasons. Both nuclei have charge $+e$ only, giving the lowest possible Coulomb
barrier. It needs no weak interaction. And it releases a great deal of energy because it forms
tightly bound helium-4. Even so it needs plasma temperatures of about \SI{e8}{K}, hotter than the Sun's
core. A useful power output requires three things at once: high temperature, high density and long
confinement time; their product must exceed a threshold called the \emph{Lawson criterion}. Magnetic
confinement devices, the tokamaks such as the international ITER project in which India is a partner,
hold the plasma away from the walls with magnetic fields. Inertial confinement crushes fuel pellets with
lasers; in December 2022 the US National Ignition Facility for the first time obtained more fusion energy
from a pellet than the laser energy delivered to it, although far less than the electricity used to run
the lasers. A power station is still decades away.

\begin{exbox}{Deuterium--tritium energetics}\label{ex:react-dt}
Using $M(\iso{2}{H})=\SI{2.014102}{\amu}$, $M(\iso{3}{H})=\SI{3.016049}{\amu}$,
$M(\iso{4}{He})=\SI{4.002603}{\amu}$ and $m_n=\SI{1.008665}{\amu}$, find $Q$ for
$\iso{2}{H}+\iso{3}{H}\to\iso{4}{He}+n$, and the kinetic energy of the neutron when the reacting nuclei
are slow.
\stepn{1} \textbf{$Q$.} $Q=[\num{2.014102}+\num{3.016049}-\num{4.002603}-\num{1.008665}]\times\SI{931.494}{\MeV}=\SI{17.6}{\MeV}$.
\stepn{2} \textbf{Sharing.} With nearly zero total momentum the products recoil back to back, and the
neutron, about a quarter of the helium's mass, takes the larger share:
$K_n=Q\,\dfrac{M_{\mathrm{He}}}{M_{\mathrm{He}}+m_n}=17.6\times\dfrac{4.003}{5.011}=\SI{14.1}{\MeV}$.
\stepn{3} \textbf{Meaning.} About four-fifths of the energy leaves with uncharged neutrons, which no
magnetic field can hold. They carry the energy out of a fusion reactor to a surrounding blanket, where it
becomes heat, and they make the reactor walls radioactive.
\ans{$Q=\SI{17.6}{\MeV}$; $K_n=\SI{14.1}{\MeV}$}
\end{exbox}

\section{Summary}\label{sec:react-summary}
A nuclear reaction is bookkept like a decay with two particles at the start: $A$ and $Z$ balance, and the
$Q$-value is the rest energy that disappears. The binding-energy curve decides the sign: nucleons moving
towards its summit near iron release energy, which is why heavy nuclei release energy by splitting and
light ones by merging.

\begin{itemize}
\item $\mathrm X(a,b)\mathrm Y$: $Q=[(m_a+m_{\mathrm X})-(m_b+m_{\mathrm Y})]c^2$ with atomic masses;
      $Q>0$ exoergic, $Q<0$ endoergic. Equivalently $Q$ is the binding energy gained.
\item Momentum conservation locks the CM kinetic energy away: the threshold of an endoergic reaction on a
      target at rest is $K_{\text{th}}=|Q|(1+m_a/m_{\mathrm X})$. For $Q>0$ with slow reactants, two
      products share $Q$ in inverse ratio of their masses.
\item Fission of uranium-235 is triggered by a slow neutron (no Coulomb barrier); one split such as
      Ba\,+\,Kr\,+\,3$n$ has $Q=\SI{173}{\MeV}$, and the average total per fission, including the
      fragments' later $\beta$ decays, is about \SI{200}{\MeV}, mostly fragment kinetic energy.
\item Multiplication factor $k$: subcritical $k<1$, critical $k=1$, supercritical $k>1$; growth $k^n$ per
      $n$ generations; a source in a subcritical assembly is multiplied by $1/(1-k)$. Leakage through the
      surface sets a critical mass.
\item A reactor holds $k=1$ with enriched fuel, a moderator, control rods and coolant. A head-on elastic
      collision leaves a neutron the fraction $\phi=((A-1)/(A+1))^2$ of its energy, so light nuclei
      moderate best; hydrogen also absorbs neutrons, so heavy water allows natural uranium.
\item Delayed neutrons ($0.65\%$) stretch the generation time from about \SI{e-4}{s} to about \SI{0.1}{s},
      provided $k-1$ stays below the delayed fraction (no prompt criticality).
\item Fusion releases about \SI{3.5}{\MeV} per nucleon (D--T) against \SI{0.85}{\MeV} for fission. The
      Coulomb barrier, $U_B=Z_1Z_2e^2/4\pi\varepsilon_0(R_1+R_2)$, is a fraction of an MeV; the classical
      temperature $2U_B/3k_B$ is billions of kelvin, but stars fuse at about \SI{1.5e7}{K} through the
      Maxwell tail and tunnelling (quoted, not derived).
\item The Sun's proton--proton chain turns $4p$ into $\iso{4}{He}$ with $Q=\SI{26.7}{\MeV}$; the Sun loses
      $P_\odot/c^2\approx\SI{4.3e9}{kg/s}$. Terrestrial fusion uses D--T ($Q=\SI{17.6}{\MeV}$, neutron
      \SI{14.1}{\MeV}) at about \SI{e8}{K}.
\end{itemize}

\section{Problem Types}\label{sec:react-types}
Seven types cover what examiners ask about reactions, fission and fusion. Nearly all reduce to two
skills: an energy balance from masses, and a momentum argument for how energy is shared or locked away.

\archq{Reaction \texorpdfstring{$Q$}{Q}-values and threshold energies}
{A reaction $\mathrm X(a,b)\mathrm Y$ or a written equation with masses; asked whether it is exoergic or
endoergic, how much energy it releases, or the minimum projectile energy.}
{Balance $A$ and $Z$ to identify any unknown particle. $Q=[(m_a+m_{\mathrm X})-(m_b+m_{\mathrm Y})]\times
\SI{931.5}{\MeV}$ per \si{\amu}, with atomic masses. If $Q<0$, $K_{\text{th}}=|Q|(1+m_a/m_{\mathrm X})$,
Eq.~\eqref{eq:react-threshold}, where $m_a$ is the \emph{moving} particle. If $Q>0$ and the projectile
is slow, the products share $Q$ in inverse ratio of their masses.}
{Taking the threshold to be $|Q|$, or putting the target's mass where the projectile's belongs.}
{The reaction $\iso{7}{Li}(p,n)\iso{7}{Be}$ has $Q=-\SI{1.645}{\MeV}$. Taking
$m_p/m_{\mathrm{Li}}=1.008/7.016$, the minimum kinetic energy of protons fired at stationary lithium-7 for
the reaction to occur is:}
{\SI{1.65}{\MeV}}
{\SI{1.44}{\MeV}}
{\SI{13.1}{\MeV}}
{\SI{1.88}{\MeV}\archans{D}{$K_{\text{th}}=1.645\,(1+1.008/7.016)=1.645\times1.144=\SI{1.88}{\MeV}$.
(A) ignores momentum conservation; (B) multiplies $|Q|$ by $m_{\mathrm X}/(m_a+m_{\mathrm X})$ instead of
dividing by it; (C) uses $1+m_{\mathrm X}/m_a$, the threshold for
lithium ions fired at stationary protons.}}

\archq{Fission energy from masses or binding energies}
{A specific fission equation with masses, or binding energies per nucleon of the parent and fragments;
energy per fission, per mole or per kilogram; comparisons with chemical fuel.}
{Balance the neutrons. Either $Q=[\sum M_{\text{in}}-\sum M_{\text{out}}]\times\SI{931.5}{\MeV}$ per
\si{\amu} (atomic masses), or $Q=\sum\Eb(\text{products})-\sum\Eb(\text{reactants})$ with
$\Eb=A\times(\Eb/A)$. Energy per kilogram $=(1000/M_{\text{molar}})\,N_A\times E_{\text{fission}}$, with
$M_{\text{molar}}$ in grams per mole.}
{Multiplying only one side's binding energy, or using the difference in $\Eb/A$ without multiplying by the
number of nucleons.}
{The excited compound nucleus $\iso{236}{U}^*$ ($\Eb/A=\SI{7.6}{\MeV}$) splits into $\iso{140}{Xe}$
($\Eb/A=\SI{8.4}{\MeV}$), $\iso{94}{Sr}$ ($\Eb/A=\SI{8.6}{\MeV}$) and two free neutrons. The energy
released is about:}
{\SI{191}{\MeV}}
{\SI{0.9}{\MeV}}
{\SI{1984}{\MeV}}
{\SI{212}{\MeV}\archans{A}{$Q=140(8.4)+94(8.6)-236(7.6)=1176+808.4-1793.6=\SI{191}{\MeV}$; the two free
neutrons have no binding energy. (B) is a typical gain per nucleon only; (C) is the products' total
binding energy, forgetting the parent's; (D) is $236\times\SI{0.9}{\MeV}$, which wrongly lets all $236$
nucleons, the free neutrons included, gain \SI{0.9}{\MeV} each.}}

\archq{Reactor power and fuel consumption}
{Power output (thermal or electrical) and fuel burnt per second, day or year.}
{Thermal power $=$ electrical power $/$ efficiency. Fissions per second $=P_{\text{th}}/E_{\text{fission}}$
with $E_{\text{fission}}$ in joules. Mass per second $=(\text{fissions per second})\times M_{\text{molar}}/N_A$.}
{Using the electrical power directly, or multiplying by the efficiency instead of dividing.}
{A reactor delivers \SI{300}{MW} of electrical power at $25\%$ efficiency. With \SI{200}{\MeV} per
fission, the number of fissions per second is about:}
{\num{9.4e18}}
{\num{3.7e19}}
{\num{2.3e18}}
{\num{3.7e25}\archans{B}{$P_{\text{th}}=300/0.25=\SI{1200}{MW}$, and
$\num{1.2e9}/(200\times\num{1.602e-13})=\SI{3.7e19}{s^{-1}}$. (A) uses the electrical power; (C)
multiplies by the efficiency; (D) converts MeV to joules with $\num{1.602e-19}$, the factor for eV.}}

\archq{Moderation by elastic collisions}
{Energy lost by neutrons striking moderator nuclei; number of collisions to thermalise.}
{Head-on elastic collision: energy kept $\phi=((A-1)/(A+1))^2$, energy transferred $1-\phi$. After $n$
collisions $K_n=K_0\phi^n$, so $n=\ln(K_n/K_0)/\ln\phi$; round up.}
{Forgetting to square: $(A-1)/(A+1)$ is the \emph{velocity} ratio; the energy ratio is its square.}
{A neutron collides head-on and elastically with a stationary deuteron (mass number $2$). The fraction
of the neutron's kinetic energy transferred to the deuteron is:}
{$1/9$}
{$1/3$}
{$8/9$}
{$2/3$\archans{C}{The neutron keeps $\phi=(1/3)^2=1/9$ and transfers $1-1/9=8/9$. (A) is the fraction
kept; (B) and (D) use the velocity ratio $1/3$ without squaring.}}

\archq{Growth of a chain reaction}
{Multiplication factor $k$, generations, generation time, power growth or decay.}
{After $n$ generations the population is $k^n$ times larger, so $n=\ln(\text{factor})/\ln k$; time
$=n\times$ generation time. For $k<1$ with a source, each source neutron gives $1/(1-k)$ neutrons in all.}
{Taking logarithms in the wrong base: $k^n=10^{24}$ needs $n=\log_k 10^{24}=\ln10^{24}/\ln k$, not
$\ln10^{24}$ or $24$.}
{In a fission assembly $k=1.5$ and the generation time is \SI{10}{ns}. Starting from one fission, the time
until a single generation contains about $10^{24}$ fissions is about:}
{\SI{0.24}{\micro s}}
{\SI{0.55}{\micro s}}
{\SI{0.80}{\micro s}}
{\SI{1.36}{\micro s}\archans{D}{$n=\ln10^{24}/\ln1.5=55.3/0.405=136$ generations, and
$136\times\SI{10}{ns}=\SI{1.36}{\micro s}$. (A) takes $n=24$ (base $10$); (B) takes $n=\ln10^{24}=55$
(base $e$); (C) takes $n=\log_210^{24}=80$, as if each fission led to two more.}}

\archq{Coulomb barrier and ignition temperature}
{Minimum energy for two nuclei to touch; temperature to supply it thermally.}
{Contact distance $r=R_1+R_2$ with $R=R_0A^{1/3}$, $R_0=\SI{1.2}{\fm}$ (or as given). $U_B=Z_1Z_2
(\SI{1.44}{\MeV\fm})/r$. Set $\tfrac32k_BT=U_B$ for the classical temperature, then note that tunnelling and
the Maxwell tail allow fusion far below it.}
{Forgetting the charge product $Z_1Z_2$, or using one radius instead of the sum of two.}
{Taking $R=1.2A^{1/3}\,\si{\fm}$, the Coulomb barrier for two $\alpha$ particles to touch is about:}
{\SI{0.38}{\MeV}}
{\SI{3.02}{\MeV}}
{\SI{1.51}{\MeV}}
{\SI{0.76}{\MeV}\archans{C}{$r=2(1.2)(4^{1/3})=\SI{3.81}{\fm}$ and $U_B=(2)(2)(1.44)/3.81=\SI{1.51}{\MeV}$.
(A) forgets the charge product $4$; (B) uses one radius, $\SI{1.91}{\fm}$; (D) uses $Z_1Z_2=2$.}}

\archq{Fusion energetics and stellar power}
{$Q$ of a fusion reaction from masses; sharing of the energy between two products; stellar mass loss and
fuel consumption.}
{$Q$ from atomic masses. Two products from nearly zero momentum share in inverse ratio of mass, exactly as
in $\alpha$ decay: $K_1=Q\,m_2/(m_1+m_2)$. Stellar mass loss $=P/c^2$; fuel rate $=(P/Q)\times$ fuel mass
per reaction.}
{Giving the larger share to the heavier product: equal momenta mean the \emph{lighter} product is faster
and carries more energy.}
{Two slow deuterons fuse: $\iso{2}{H}+\iso{2}{H}\to\iso{3}{He}+n$, with $Q=\SI{3.27}{\MeV}$. Taking the
helium-3 mass as three times the neutron's, the kinetic energy of the neutron is about:}
{\SI{2.45}{\MeV}}
{\SI{0.82}{\MeV}}
{\SI{1.64}{\MeV}}
{\SI{3.27}{\MeV}\archans{A}{$K_n=Q\times\frac{3}{3+1}=\frac34(3.27)=\SI{2.45}{\MeV}$. (B) is the
helium-3 share; (C) splits equally; (D) gives everything to the neutron, ignoring the recoil that momentum
conservation forces.}}

\section{Practice Problems}\label{sec:react-practice}

The problems below are of standard examination difficulty and are meant to be worked without looking back
at the chapter. Answers, with brief working, follow the list. Use these data throughout:
\begin{center}\small
\begin{tabular}{@{}lll@{}}
$1\,\si{\amu}\,c^2=\SI{931.494}{\MeV}$ & $m_n=\SI{1.008665}{\amu}$ & $M(\iso{1}{H})=\SI{1.007825}{\amu}$\\
$M(\iso{2}{H})=\SI{2.014102}{\amu}$ & $M(\iso{3}{H})=\SI{3.016049}{\amu}$ & $M(\iso{3}{He})=\SI{3.016029}{\amu}$\\
$M(\iso{4}{He})=\SI{4.002603}{\amu}$ & $M(\iso{235}{U})=\SI{235.043930}{\amu}$ & $N_A=\SI{6.022e23}{mol^{-1}}$\\
$k_B=\SI{1.381e-23}{J/K}$ & $c=\SI{2.998e8}{m/s}$ & $\SI{1}{\MeV}=\SI{1.602e-13}{J}$
\end{tabular}
\end{center}


\begin{enumerate}[itemsep=3pt]
\item Find $Q$ for Chadwick's reaction $\iso{9}{Be}(\alpha,n)\iso{12}{C}$. The atomic masses are
      \SI{9.012183}{\amu} for beryllium-9 and \SI{12.000000}{\amu} for carbon-12. Is the reaction
      exoergic or endoergic?
\item Cobalt-60 is made by neutron capture, $\iso{59}{Co}(n,\gamma)\iso{60}{Co}$. The atomic masses
      are \SI{58.933194}{\amu} for cobalt-59 and \SI{59.933816}{\amu} for cobalt-60. Find the energy
      carried by the gamma radiation (neglect the recoil of the nucleus).
\item Rutherford's reaction $\iso{14}{N}(\alpha,p)\iso{17}{O}$ has $Q=-\SI{1.19}{\MeV}$. Explain why an
      $\alpha$ particle of exactly \SI{1.19}{\MeV} cannot make it go, and find the true threshold (take
      the masses as $4.003$ and $14.003\,\si{\amu}$). Would the threshold be higher or lower if, instead,
      nitrogen ions were fired at a stationary helium target?
\item Another fission channel is $n+\iso{235}{U}\to\iso{140}{Xe}+\iso{94}{Sr}+2n$. Given
      $M(\iso{140}{Xe})=\SI{139.921646}{\amu}$ and $M(\iso{94}{Sr})=\SI{93.915356}{\amu}$, find $Q$.
\item The Hiroshima bomb released energy equivalent to about $15$ kilotonnes of TNT
      ($1$~kilotonne $=\SI{4.184e12}{J}$). Taking \SI{200}{\MeV} per fission, estimate the mass of
      uranium-235 that actually fissioned.
\item The fragments $\iso{141}{Ba}$ ($Z=56$) and $\iso{92}{Kr}$ ($Z=36$) are produced in fission. The
      heaviest stable barium isotope is $\iso{138}{Ba}$ and the heaviest stable krypton isotope is
      $\iso{86}{Kr}$. Compare the neutron numbers $A-Z$ and state what kind of decays follow.
\item A plant delivers \SI{1000}{MW} of \emph{electrical} power at $33.3\%$ efficiency. How much
      uranium-235 does it fission per day and per year (\SI{200}{\MeV} per fission)?
\item A coal-fired plant producing \SI{1000}{MW} of \emph{thermal} power burns coal of heat value
      \SI{30}{MJ/kg}. How many kilograms of coal does it need per day? Compare with the \SI{1.05}{kg} of
      uranium-235 per day that a reactor of the same thermal power fissions.
\item Beryllium-9 is used as a moderator and reflector. How many head-on elastic collisions with
      beryllium-9 nuclei slow a \SI{2.0}{\MeV} neutron to \SI{0.025}{eV}?
\item In a single head-on collision with a proton, a neutron stops. Why, then, is ordinary water a less
      economical moderator than heavy water?
\item Derive Eq.~\eqref{eq:react-phi} from conservation of momentum and kinetic energy.
\item In a fission weapon $k\approx2$ and the generation time is about \SI{10}{ns}. Starting from one
      fission, how many fissions occur in the 80th generation alone, how long does it take to reach that
      generation, and what energy (in kilotonnes of TNT) do those fissions release at \SI{200}{\MeV} each?
      Roughly how much do all eighty generations release together?
\item A subcritical assembly with $k=0.90$ contains a neutron source. Counting the source neutron itself,
      how many neutrons in total does each source neutron eventually give rise to? What happens to this
      number as $k\to1$?
\item\label{prob:react-pcbar} In stars heavier than the Sun, hydrogen also burns by proton capture on
      carbon-12 (the CNO cycle). Taking $R=R_0A^{1/3}$ for both nuclei, estimate the Coulomb barrier
      between a proton and a carbon-12 nucleus, and the temperature at which $\tfrac32k_BT$ equals it.
\item At the Sun's centre (\SI{1.5e7}{K}), $k_BT=\SI{1.29}{\keV}$. By what factor is it below the
      proton--carbon barrier of Problem~\ref{prob:react-pcbar}, and below the proton--proton barrier of
      Section~\ref{sec:react-fusion}? Which reaction would you expect to need the hotter star?
\item Why must stars reach about \SI{e8}{K} before helium can fuse, much hotter than hydrogen burning
      requires?
\item Two deuterons can fuse by either of two branches: $\iso{2}{H}+\iso{2}{H}\to\iso{3}{He}+n$ or
      $\iso{2}{H}+\iso{2}{H}\to\iso{3}{H}+p$. Find $Q$ for each.
\item The Sun loses mass at \SI{4.3e9}{kg/s} and converts \SI{6.0e11}{kg} of hydrogen into helium
      every second (Section~\ref{sec:react-sun}). (a) What fraction of its mass (\SI{2.0e30}{kg}) has it
      lost in its $4.6$ billion years? (b) Only the core, roughly $10\%$ of the Sun's hydrogen, is hot
      enough to burn. Taking the Sun's mass to be $71\%$ hydrogen, estimate its hydrogen-burning
      lifetime at today's rate.
\item Compare the energy released per kilogram of deuterium--tritium fuel ($Q=\SI{17.6}{\MeV}$) with that
      per kilogram of uranium-235 (\SI{200}{\MeV} per fission).
\end{enumerate}

\subsection*{Answers}
{\small
\begin{enumerate}[itemsep=2pt]
\item $Q=[\num{9.012183}+\num{4.002603}-\num{12.000000}-\num{1.008665}]\times\SI{931.494}{\MeV}=+\SI{5.70}{\MeV}$:
      exoergic.
\item $Q=[\num{58.933194}+\num{1.008665}-\num{59.933816}]\times\SI{931.494}{\MeV}=\SI{7.49}{\MeV}$, emitted
      as gamma rays: it is the binding energy of the captured neutron in cobalt-60.
\item At exactly \SI{1.19}{\MeV} the products would have to be at rest, with zero momentum, but the
      $\alpha$ particle brought in momentum; part of its energy must stay as CM motion.
      $K_{\text{th}}=1.19(1+4.003/14.003)=\SI{1.53}{\MeV}$. Firing nitrogen at helium makes the projectile
      the heavy partner: $K_{\text{th}}=1.19(1+14.003/4.003)=\SI{5.36}{\MeV}$, much higher, because now
      most of the projectile's energy is locked in the motion of the CM.
\item $Q=[\num{235.043930}+\num{1.008665}-\num{139.921646}-\num{93.915356}-2(\num{1.008665})]\times\SI{931.494}{\MeV}=\SI{185}{\MeV}$.
\item $E=15\times\num{4.184e12}=\SI{6.3e13}{J}$; at \SI{8.2e13}{J} per kilogram (Example~\ref{ex:react-perkg}),
      $m=\num{6.3e13}/\num{8.2e13}\approx\SI{0.76}{kg}$, out of about \SI{64}{kg} of uranium in the weapon.
\item Barium-141 has $A-Z=85$ against $82$ for the heaviest stable barium; krypton-92 has $A-Z=56$ against
      $50$ for the heaviest stable krypton. Both fragments have too many neutrons (they lie above the
      valley of stability), so each undergoes a chain of $\beta^-$ decays towards it, for example
      $\iso{141}{Ba}\to\iso{141}{La}\to\iso{141}{Ce}\to\iso{141}{Pr}$ (stable).
\item Thermal power $1000/0.333=\SI{3000}{MW}$, three times Example~\ref{ex:react-fuel}: $3\times1.05=\SI{3.16}{kg}$
      per day, about \SI{1150}{kg} per year.
\item $\num{1e9}\times\num{86400}/\num{30e6}=\SI{2.9e6}{kg}$ (about $2900$ tonnes) per day, some $2.7$
      million times the \SI{1.05}{kg} of uranium-235.
\item $\phi=(8/10)^2=0.64$; $n=\ln(\num{1.25e-8})/\ln0.64=(-18.2)/(-0.446)=40.8$, so $41$ collisions,
      between deuterium's $9$ and carbon's $55$.
\item Hydrogen nuclei also absorb neutrons ($n+p\to\iso{2}{H}+\gamma$), wasting them. Light-water
      reactors must compensate with enriched fuel; heavy water absorbs far fewer neutrons and allows
      natural uranium.
\item Momentum gives $v_0-v_1=AV$; energy gives $(v_0-v_1)(v_0+v_1)=AV^2$; dividing, $v_0+v_1=V$. Then
      $v_0-v_1=A(v_0+v_1)$, so $v_1=-v_0(A-1)/(A+1)$ and $\phi=v_1^2/v_0^2=((A-1)/(A+1))^2$.
\item $2^{80}=\num{1.2e24}$ fissions in the 80th generation, reached after $80\times\SI{10}{ns}=\SI{0.8}{\micro s}$;
      $E=\num{1.2e24}\times\SI{3.2e-11}{J}=\SI{3.9e13}{J}$, about \SI{9}{kt}. All eighty generations
      together hold $1+2+\dots+2^{80}=2^{81}-1$, twice as many, so about \SI{18}{kt}. (A real weapon blows
      itself apart after about this many generations, which is why only a kilogram or so of its uranium
      fissions.)
\item $1+k+k^2+\dots=1/(1-k)=10$ neutrons. As $k\to1$ the total diverges: the assembly becomes
      self-sustaining.
\item $r=1.2(1+12^{1/3})=\SI{3.95}{\fm}$; $U_B=6(1.44)/3.95=\SI{2.19}{\MeV}$;
      $T=2U_B/3k_B=\SI{1.7e10}{K}$, three times the proton--proton estimate.
\item $2190/1.29\approx1700$ for proton--carbon, against $720/1.29\approx560$ for proton--proton. The
      higher, wider barrier makes tunnelling far less likely, so the CNO cycle needs a hotter core; it
      dominates only in stars somewhat heavier than the Sun.
\item The barrier $Z_1Z_2e^2/4\pi\varepsilon_0r$ grows with the charge product, $2\times2=4$ for two
      $\alpha$ particles against $1$ for two protons; even allowing for the larger radii, the barrier is
      about twice as high (\SI{1.5}{\MeV} against \SI{0.7}{\MeV}). Tunnelling falls off exponentially with
      the charge product and with the mass of the tunnelling particles, so the heavier, doubly charged
      $\alpha$ particles fuse at a useful rate only at a far higher temperature.
\item $Q_1=[2(\num{2.014102})-\num{3.016029}-\num{1.008665}]\times\SI{931.494}{\MeV}=\SI{3.27}{\MeV}$;
      $Q_2=[2(\num{2.014102})-\num{3.016049}-\num{1.007825}]\times\SI{931.494}{\MeV}=\SI{4.03}{\MeV}$.
\item (a) $\num{4.3e9}\times4.6\times10^{9}\times\num{3.156e7}=\SI{6.2e26}{kg}$, only $\num{3e-4}$ of its
      mass. (b) Available hydrogen $0.10\times0.71\times\num{2.0e30}=\SI{1.4e29}{kg}$; at
      \SI{6.0e11}{kg/s} it lasts $\SI{2.4e17}{s}\approx\SI{7.5e9}{\year}$, the right order for the Sun's
      ten-billion-year life.
\item $\SI{17.6}{\MeV}$ per $\SI{5.03}{\amu}$ of fuel gives $\SI{3.4e14}{J/kg}$, about $4.1$ times
      uranium-235's \SI{8.2e13}{J/kg}.
\end{enumerate}
}
